\chapter{Introduction to Discrete Mathematics}
\signature

\section{What is discrete mathematics?}

Discrete mathematics is the study of mathematical structures that are
fundamentally \emph{discrete} rather than continuous: integers, finite sets,
graphs, logical statements, sequences and algorithms. Where calculus deals with
quantities that vary smoothly, discrete mathematics deals with objects that can
be counted, listed, and manipulated one at a time.

\begin{definition}
\emph{Discrete mathematics} is the branch of mathematics dealing with countable
structures --- objects that can assume only distinct, separated values --- as
opposed to continuous structures such as the real line.
\end{definition}

The natural habitat of discrete mathematics is the world of integers
$\Z = \{\dots,-2,-1,0,1,2,\dots\}$ and of finite collections. A digital
computer is itself a discrete object: its memory holds finitely many bits, each
either $0$ or $1$, so every question about computation is ultimately a question
of discrete mathematics.

\section{The main questions}

Across all its subfields, discrete mathematics keeps returning to a small
number of fundamental questions.

\begin{itemize}
\item \textbf{Existence.} Is there an object with a given property? (Is there a
  path visiting every vertex of a graph exactly once?)
\item \textbf{Counting.} How many such objects are there? (How many ways can a
  committee of $5$ be chosen from $30$ students?)
\item \textbf{Optimization.} Which object is best? (Which spanning tree of a
  network has minimal total cost?)
\item \textbf{Structure.} How are the objects organised? (How does a partial
  order arrange the subsets of a set?)
\item \textbf{Algorithms.} How can the answer be computed, and how fast? (Can
  satisfiability of a logical formula be decided efficiently?)
\end{itemize}

\section{Information, summarization and computation}

Three recurring activities motivate the tools we will build.

\subsection{Information}
Discrete structures \emph{encode information}. A truth table encodes the
behaviour of a logical circuit; a set encodes membership; a relation encodes
links between objects; a binary string encodes a file. Understanding which
structures can encode which information --- and how compactly --- is a central
concern.

\subsection{Summarization}
Mathematics compresses infinitely many facts into a single statement. The
formula
\[
\sum_{i=1}^{n} i \;=\; \frac{n(n+1)}{2}
\]
summarises infinitely many arithmetic identities at once. Quantifiers
($\forall$, $\exists$), closed forms and recurrences are all instruments of
summarization, and proving such summaries correct is what Part~IV of this book
is about.

\subsection{Computation}
Every algorithm is a discrete process: a sequence of elementary steps acting on
discrete data. Counting the steps (complexity), proving the result correct
(invariants, induction) and even defining what ``computable'' means are tasks
for discrete mathematics.

\section{A brief history}

Discrete themes are as old as mathematics itself. Euclid (c.~300 BCE) proved
that there are infinitely many primes --- a discrete existence theorem we shall
reprove in Chapter~12. Combinatorial counting flourished with Pascal and
Fermat in the 17th century around games of chance, producing Pascal's triangle
(Chapters~10--11). Euler's 1736 solution of the K\"onigsberg bridge problem
founded graph theory. In the 19th century, Boole turned logic into algebra and
Cantor created set theory, giving the subject its modern foundations
(Parts~I--II). The 20th century brought the decisive turn: G\"odel, Turing and
Church connected logic to computation, and the arrival of digital computers
transformed discrete mathematics from a collection of curiosities into the
mathematical backbone of computer science.

\section{Discrete mathematics today}

Modern applications are everywhere:

\begin{itemize}
\item \textbf{Computer science} --- data structures, databases, programming
  language semantics, verification of software and hardware.
\item \textbf{Cryptography} --- number theory and finite algebra protect every
  online transaction.
\item \textbf{Networks} --- graph theory models the internet, social networks
  and transportation systems.
\item \textbf{Data science} --- counting, probability on finite spaces, and
  combinatorial optimisation underlie machine learning pipelines.
\item \textbf{Operations research} --- scheduling, allocation and routing are
  discrete optimisation problems.
\end{itemize}

\begin{notebox}
Continuous mathematics asks ``how much?''; discrete mathematics asks
``how many?'', ``does one exist?'', and ``how can we build it?''.
\end{notebox}

\section{Why study discrete mathematics?}

Beyond its applications, discrete mathematics teaches a craft: the craft of
\emph{precise reasoning}. In this course you will learn to

\begin{itemize}
\item read and write statements whose meaning is completely unambiguous
  (Chapters~2--5);
\item manipulate the basic structures of all mathematics --- sets, relations,
  functions, sequences (Chapters~6--9);
\item count complicated configurations exactly (Chapters~10--11);
\item construct proofs that withstand any scrutiny, including proofs by
  induction and recursive constructions (Chapters~12--14).
\end{itemize}

These skills transfer directly to programming, to algorithm design and to every
mathematical subject you will meet later.

\section{About this course}

The book is organised in four parts.

\begin{itemize}
\item \textbf{Part I --- Foundations of Logic.} Propositions, connectives,
  truth tables, logical equivalence, predicates, quantifiers and rules of
  inference: the language in which all later chapters are written.
\item \textbf{Part II --- Discrete Structures.} Sets and their operations,
  relations, functions, sequences and series: the objects that the language
  talks about.
\item \textbf{Part III --- Counting.} Permutations, combinations, the binomial
  theorem and a toolbox of binomial identities.
\item \textbf{Part IV --- Proofs and Recursion.} Proof techniques, mathematical
  induction in its weak, strong and structural forms, and recursive
  definitions of functions, sets and structures.
\end{itemize}

Each chapter ends with a graded set of exercises --- tagged
\easy{}, \medium{} or \hard{} --- followed by complete solutions.

\section{Who is this course for?}

The course assumes no prerequisites beyond secondary-school algebra. It is
designed for first-year students of computer science, mathematics, engineering
and data science, and for anyone who wants to reason precisely about discrete
problems. Mathematical maturity is not required at the start --- building it is
precisely the point.

\section*{Exercises}
\addcontentsline{toc}{section}{Exercises}

\begin{exercise}{\easy}
Classify each of the following as a \emph{discrete} or a \emph{continuous}
question: (a) How many $5$-card poker hands are possible? (b) What is the area
under the curve $y=x^2$ between $0$ and $1$? (c) Is there a way to schedule
$10$ exams in $4$ slots with no student conflict? (d) At what speed does a
falling object hit the ground?
\end{exercise}

\begin{exercise}{\easy}
Give one example from everyday life of each of the five main questions of
discrete mathematics (existence, counting, optimization, structure,
algorithms).
\end{exercise}

\begin{exercise}{\medium}
A byte consists of $8$ bits, each $0$ or $1$. How many distinct bytes are
there? More generally, how many distinct bit strings of length $n$ are there?
Justify your count.
\end{exercise}

\begin{exercise}{\medium}
The formula $\sum_{i=1}^{n} i = \frac{n(n+1)}{2}$ ``summarises infinitely many
facts''. Verify it directly for $n = 1, 2, 3, 4$, and explain why no amount of
direct verification constitutes a proof for all $n$.
\end{exercise}

\begin{exercise}{\hard}
Euclid proved there are infinitely many primes. Suppose, for contradiction,
that $p_1, p_2, \ldots, p_k$ were \emph{all} the primes, and consider
$N = p_1 p_2 \cdots p_k + 1$. Show that none of the $p_i$ divides $N$, and
deduce a contradiction. (You may use the fact that every integer greater than
$1$ has at least one prime divisor.)
\end{exercise}

\section*{Solutions to the Exercises}
\addcontentsline{toc}{section}{Solutions to the Exercises}

\begin{solution}{1.1}
(a) Discrete --- a finite count. (b) Continuous --- an integral over a real
interval. (c) Discrete --- existence of a finite arrangement. (d) Continuous
--- a real-valued physical quantity varying smoothly with time.
\end{solution}

\begin{solution}{1.2}
For instance: \emph{existence} --- is there a seating plan where no two rivals
sit together? \emph{Counting} --- how many phone numbers are possible with a
given prefix? \emph{Optimization} --- which route to work is shortest?
\emph{Structure} --- how is a family tree organised? \emph{Algorithms} --- what
procedure sorts a list of names alphabetically? (Any sensible examples are
acceptable.)
\end{solution}

\begin{solution}{1.3}
Each of the $8$ bits can be chosen independently in $2$ ways, so there are
$2 \cdot 2 \cdots 2 = 2^8 = 256$ bytes. The same argument gives $2^n$ bit
strings of length $n$: choosing a string means making $n$ independent binary
choices, and by the product rule (Chapter~10) the total is $2^n$.
\end{solution}

\begin{solution}{1.4}
$n=1$: $1 = \frac{1\cdot 2}{2}$. $n=2$: $1+2 = 3 = \frac{2\cdot 3}{2}$.
$n=3$: $1+2+3=6=\frac{3\cdot4}{2}$. $n=4$: $1+2+3+4=10=\frac{4\cdot5}{2}$.
Direct verification only checks finitely many instances, while the claim is
about \emph{all} natural numbers $n$; no finite list of checks can rule out a
failure at some larger $n$. A general argument --- such as induction
(Chapter~13) or Gauss's pairing trick --- is required.
\end{solution}

\begin{solution}{1.5}
Suppose $p_1,\dots,p_k$ are all the primes and let
$N = p_1p_2\cdots p_k + 1$. For each $i$, dividing $N$ by $p_i$ leaves
remainder $1$, since $p_i \mid p_1\cdots p_k$; hence $p_i \nmid N$. But
$N > 1$, so $N$ has some prime divisor $q$. Then $q$ is a prime different from
every $p_i$, contradicting the assumption that the list contained all primes.
Therefore there are infinitely many primes.
\end{solution}
